\section{The exact cost of blocking a torsion search}
\label{sec:threshold}
In the Lean proof, \texttt{exists\_rootOfUnity\_eval\_le\_gaussVal\_add}
compares a polynomial at a primitive $p$-power root minus one with its Gauss
valuation at that root's radius, with an additive tolerance. It selects a
shell using a harmonic-sum bound and controls contributions from polynomial
roots at that radius. The later simultaneous selection theorem controls
individual polynomials near their respective witness values; the Selmer
argument uses the construction to choose twists. We take up the fixed-radius
variant in the incoming reports.
\ghsource{grh/src/New/Congruences/Selmer/RootsOfUnity.lean\#L292-L305}{Single polynomial};
\ghsource{grh/src/New/Congruences/Selmer/RootsOfUnity.lean\#L474-L486}{simultaneous selection};
\ghsource{grh/src/New/Congruences/Selmer/GeneralTheorem.lean\#L318-L344}{Selmer application}.

One prior report asks for the threshold on a prescribed interval of torsion
orders over a ramified base, including intervals which cannot be
specified by a valuation cutoff. It also asks for the sharp joint dependence
on degree, tolerance and order interval. We answer both questions for
a single polynomial with unrestricted integral coefficients. The input can
in fact be any finite set of points in the open unit disc.
\href{https://github.com/VladimirReshetnikov/ProveIt/blob/d31039a48d963692e5a5176f47b6603fc9e9d9cf/docs/incoming/grh_method_transfers_sharp_torsion.zip}{Prior report, section on formal height and final questions}.

Normalize $v(p)=1$ on $\mathbb C_p$, put $v(0)=+\infty$, and fix a
nonempty finite $S\subset\{z:v(z)>0\}$, a real $\rho>0$ and a tolerance
$\delta\ge0$. The coefficient ring is $\mathcal O_{\mathbb C_p}$.
For a nonzero integral polynomial define its Gauss valuation by
\[
 A(g)=G_\rho(g)=\min_i(v(a_i)+i\rho),\qquad g(T)=\sum_i a_iT^i.
\]
A polynomial \emph{blocks} $S$ when $v(g(z))>A(g)+\delta$ for every $z\in S$.
A point attaining equality is a successful test point.
We compute the smallest possible $A(g)$ of a blocker and its least possible degree.

Partition $S$ by the positive weights $w(z)=\min(\rho,v(z))$ and process
the distinct weights $w$ in decreasing order. Denote the part at $w$ by
$S_w$. On each part take a greedy sequence $z_0,\ldots,z_{l-1}$, $l=|S_w|$:
start at any point, and at step $m<l$ minimize the following quantity over
the unused points $x$:
\[
 E_w(m)=\min_x\sum_{i<m}\bigl(v(x-z_i)-w\bigr),\qquad E_w(l)=+\infty.
\]
In particular $E_w(0)=0$. Given the previously computed integers $m_u$, put
\begin{equation}\label{eq:rec}
 t_w=\delta+\sum_{u>w}m_u(u-w),\qquad
 m_w=\min\{m:E_w(m)>t_w\}.
\end{equation}
The computation uses only the distances in each part and its weight.

\begin{theorem}\label{th:exact}
The least blocking degree $d_*$ and least blocking budget $B_*$ are
\[
 d_*=\sum_w m_w,\qquad B_*=\sum_w w m_w.
\]
One monic polynomial attains both: take the first $m_w$ greedy points of each
part as its roots. Consequently a polynomial with $\deg g\le d$ and
$G_\rho(g)\le B$ is guaranteed a successful test point precisely when
$d<d_*$ or $B<B_*$. With no degree restriction the budget threshold alone
also applies to integral formal power series.
\end{theorem}

\begin{proof}
First consider one part $U=S_w$. Write $P_m=\prod_{i<m}(T-z_i)$ and
$b_m=mw+E_w(m)$ for $0\le m<|U|$. Greedy choice says
$\min_U v(P_m)=v(P_m(z_m))=b_m$. Among all monic polynomials of degree
$m$, this is the largest possible minimum. Indeed write a competing
monic $Q=\sum_{i=0}^m c_i P_i$, with $c_m=1$. If $v(Q(x))>b_m$ for all
$x\in U$, evaluating successively at $z_0,\ldots,z_m$ forces
$v(c_i)+b_i>b_m$ for every $i$. All earlier terms exceed the threshold
everywhere on $U$ and later terms vanish at $z_i$. At $i=m$ we have a
contradiction. Thus tie choices give the same value $E_w(m)$.
At $m=|U|$ the root polynomial vanishes on the whole part.
This is the finite-set argument for Bhargava's $P$-orderings.
\href{https://doi.org/10.1006/jnth.1998.2220}{Bhargava on fixed divisors};
\href{https://arxiv.org/abs/1611.02142}{Amini on tree factorials}.

To see why each part is charged its stated weight, factor a prospective
blocker $g=b\prod_\alpha(T-\alpha)$ in $\mathbb C_p$. Multiplicativity
of the Gauss valuation gives
\[
 G_\rho(g)-G_0(g)=\sum_\alpha
 [\min(\rho,v(\alpha))-\min(0,v(\alpha))],\qquad G_0(g)\ge0.
\]
The contribution of this root to $v(g(z))-G_\rho(g)$ is
$v(z-\alpha)-\min(\rho,v(\alpha))$. Roots with $v(\alpha)\le0$
contribute zero. For any remaining root put $u=\min(\rho,v(\alpha))$.
If there is no part at $u$, the root contributes at most zero at every
test point and can be discarded. If the part exists, its contribution
at every other part of weight $w$ is $\min(u,w)-u$. Within $S_u$ it
comes from the monic product of all roots at weight $u$, followed by
subtraction of $u$ for each root. This reasoning includes the part at
$u=\rho$, which can contain several actual radii or the origin.

Let $M_u$ be the count of retained roots of weight $u$. At each point
of $S_w$ blocking requires the monic product of its $M_w$ roots to have
excess over $M_ww$ strictly greater than
$\delta+\sum_{u>w}M_u(u-w)$. Discarded terms could only have lowered
the excess. By the one-part extremum and induction from the largest
weight down, $M_w\ge m_w$ at every part (the claim is automatic for
$M_w\ge |S_w|$). All retained roots cost their weight in
$G_\rho(g)-G_0(g)$, proving both lower bounds. For the stated monic
polynomial the cross-part effects are exactly the charges in
\eqref{eq:rec}, the one-part minima are attained, and $G_0=0$.
This proves sharpness including both boundary equalities.

For an integral formal series with $G_\rho=A<\infty$, truncate at degree
$D$ with $(D+1)\min_{z\in S}w(z)>A+\delta$. The discarded tail
has valuation above $A+\delta$ at each point and leaves the Gauss
valuation unchanged. The attaining polynomial supplies the reverse direction.
\end{proof}

Updating one running sum per unused point evaluates the greedy sequences
in $O(\sum_w |S_w|^2)$ valuation additions/comparisons given the distances.
The polynomial may require a finite extension of the base field when the
test points are algebraic. Prescribing the field of its coefficients adds
Galois constraints and can change the answer.
For cyclotomic shells of one valuation, earlier reports
\href{https://github.com/VladimirReshetnikov/ProveIt/blob/d31039a48d963692e5a5176f47b6603fc9e9d9cf/docs/incoming/grh_sharp_torsion_and_symplectic_blocks.zip}{already computed a shell recursion}
in terms of the regular cyclotomic tree. Theorem~\ref{th:exact} uses
the actual metric on each selected part, with arbitrary clipping at $\rho$.

\subsection{Obtaining the metric from torsion counts}
Let $F$ be a one-dimensional commutative formal group of finite
reduction height over the integers of a finite extension of $\mathbb Q_p$.
For an interval $h_0\le h\le H$, take $G=F[p^H]$ in the open unit disc and
$D=F[p^{h_0-1}]$. The target is exactly $S=G\setminus D$.
The group law gives $v(x-y)=v(x\mathbin{\ominus_F}y)$: formal subtraction
is divisible by $X-Y$ with unit quotient. Hence $V_r=\{x\in G:v(x)\ge r\}$
is a subgroup and its cosets give the balls at level $r$.

For each distinct torsion valuation, the counts $|V_r|$ and
$|D\cap V_r|$ come from the slopes, with multiplicities, of the
Weierstrass polynomials of $[p^H]$ and $[p^{h_0-1}]$, including the origin
in each subgroup. At the last level use $V_\infty=\{0\}$.
At adjacent levels $r<t$, a coset of $V_r$ has $|V_r|/|V_t|$ children.
If it meets $D$, exactly $|D\cap V_r|/|D\cap V_t|$ children meet $D$,
by translation from a point of intersection. If it misses $D$, all its
children miss $D$. Also distinguish the unique path containing the origin.
This constructs the finite tree up to permutations of equivalent children.
Delete the $D$ leaves. The last common ball of two surviving leaves gives
the valuation of their difference. The analogous comparison with the
marked origin gives each weight. Theorem~\ref{th:exact} now solves the
prescribed order interval using integer populations and rational slopes;
there is no need to solve individually for the torsion coordinates.
For a fixed lower endpoint the first guaranteed upper endpoint, whenever
it exists, is the first $H$ whose computed threshold passes either strict test.
For computation, the input consists of the indicated Newton polygons or
a formal group specified effectively enough to compute them.
\href{https://arxiv.org/abs/math/0604252}{Fargues on ramification filtrations};
\href{https://arxiv.org/abs/0910.3323}{Rabinoff on canonical subgroups}.

For each fixed $B$, increasing $H$ eventually guarantees success even at tolerance zero.
Choose $0<t\le\rho$ smaller than all nonzero valuations in $D$,
with this last condition vacuous for $D=\{0\}$. All torsion points
of valuation $>t$ form a finite subgroup $N$. Indeed,
$[p](X)=pX+O(X^2)$ with integral coefficients raises any valuation
$v>t$ by at least $\min(1,t)$ until the point is killed. A nonzero
point of valuation $>1$ cannot be torsion since from then on the
linear term always uniquely dominates. Thus $N$ is killed by a
bounded power of $p$, and $D\subseteq N\subseteq G$ for large $H$.
Take one representative $x$ of each nonzero coset in $G/N$. With
$c=\sum_{0\ne z\in N}v(z)$ the sum of their valuations is
$C=(H-c)/|N|$: the sum over all nonzero torsion points in $G$ is
$H$, by the derivative $p^H$ at zero and Weierstrass preparation.
The product polynomial $P$ of the representative linear factors
has $G_t(P)=C$ and $v(P'(x))=C-v(x)$, by translation in $G/N$.
All Lagrange polynomials on these nodes have $G_t=0$. If $G_t(g)<C$,
monic division by $P$ leaves an integral remainder with Gauss value
$G_t(g)$, and interpolation gives some node with
$v(g(x))\le G_t(g)\le G_\rho(g)$. Taking $C>B$ proves termination and
explains why the preceding finite algorithm determines the optimal
endpoint as well as testing each prescribed window. This argument uses
the valuation-cutoff interpolation from the prior formal-height report.
